\chapter{Design}
\label{ch:design}

This chapter translates the requirements and architectural overview established in
Chapter~\ref{ch:analysis} into concrete technical decisions. The domain model defines the
entities and their relationships. The technology analysis justifies each major component of the
stack by comparing alternatives against the project requirements. The database schema and
entity-relationship diagram document the persistent data model. The finalised system architecture
and API design are presented, and the chapter closes with the user interface design.

% ----------------------------------------------------------------
\section{Domain Model}
\label{sec:domain-model}

The domain model consists of the entities that carry the application's persistent state. Each
entity is described below, and the complete UML class diagram is presented in
Figure~\ref{fig:uml-class-diagram} at the end of this section.

\textbf{User} is the central actor entity. It stores the account credentials (username, email,
and bcrypt-hashed password), the display name, a profile picture URL, and a role field that
takes one of two values: \texttt{STUDENT} or \texttt{INSTRUCTOR}. The entity also tracks a
cumulative study-time counter and a current-streak counter that are updated by the gamification
subsystem. A user can create many \texttt{Course} records (as an Instructor), own many
\texttt{StudentWorkspace} records (as a Student), and participate in many \texttt{StudySession}
records through the \texttt{SessionParticipant} junction entity.

\textbf{Course} represents a published course created by one Instructor. It carries a title, a
text description, a boolean \texttt{isPublished} flag that controls student visibility, and a
many-to-one relationship to its owning \texttt{User} (the Instructor). Deleting a
\texttt{Course} cascades to all its \texttt{CourseSection} and \texttt{CourseEnrollment} records
via \texttt{CascadeType.ALL} with orphan removal, ensuring no orphaned sections or enrollments
remain after a course is removed.

\textbf{CourseSection} organises Course Materials within a course into named, ordered groups. It
carries a title, an optional text description, and an \texttt{orderIndex} integer that determines
the display sequence. Each section belongs to exactly one \texttt{Course}, and each section owns
an ordered list of \texttt{CourseMaterial} records. Deleting a section cascades to all its
materials.

\textbf{CourseMaterial} represents a single file uploaded into a course section. It stores a
display title, the file storage URL, the inferred file type as a \texttt{MaterialType} enum
(\texttt{PDF}, \texttt{SLIDES}, \texttt{VIDEO}, \texttt{IMAGE}, \texttt{OTHER}), the original
filename, an optional \texttt{contributorName} field (populated when the material originated
from an approved Merge Proposal), and a timestamp. Each \texttt{CourseMaterial} belongs to
exactly one \texttt{CourseSection}.

\textbf{CourseEnrollment} is the join entity that records a student's membership in a course.
It carries an \texttt{EnrollmentStatus} field (\texttt{ACTIVE} or \texttt{DROPPED}) and the
timestamp of the initial enrollment. A unique constraint on the pair \texttt{(course\_id,
student\_id)} prevents duplicate enrollment rows. Setting the status to \texttt{DROPPED} rather
than deleting the row preserves the enrollment history and allows re-activation.

\textbf{StudentWorkspace} is the root entity of a student's Private Workspace. It holds a name,
a description, a reference to the owning \texttt{User}, and a list of \texttt{WorkspaceSpace}
records. The workspace is the container that the student manages independently; it does not belong
to any course and cannot be seen or modified by the Instructor.

\textbf{WorkspaceSpace} represents a single cloned or manually created space within a Private
Workspace. It carries a title, a description, an optional foreign key to the \texttt{Course} it
was forked from (\texttt{forkedFromCourse}), and a boolean \texttt{isPublished} flag used when a
space is surfaced in a public gallery. Each space owns an ordered list of
\texttt{WorkspaceSection} records.

\textbf{WorkspaceMaterial} mirrors \texttt{CourseMaterial} within the student's space. Two flags
encode the copy-on-write behaviour: \texttt{isReference = true} means the material points to
the original course file URL and no physical copy was made; \texttt{isHidden = true} implements
a soft-delete that hides the row from the student's view without physically removing it, which is
necessary because a contribution proposal may hold a foreign-key reference to the material row
and deleting it would violate referential integrity.

\textbf{ContributionProposal} links a student's \texttt{WorkspaceMaterial} to a target
\texttt{Course} and optionally a target \texttt{CourseSection}. Its \texttt{ProposalStatus} enum
field (\texttt{PENDING}, \texttt{APPROVED}, \texttt{REJECTED}) drives the Instructor review
workflow. When no existing section is targeted, the \texttt{proposedSectionTitle} field holds the
name of a new section to be created on approval. The \texttt{reviewedAt} timestamp is set when
the Instructor takes action, and the \texttt{contributorDisplayName} is copied from the student's
profile at submission time so it remains stable even if the student later changes their name.

\textbf{Conversation} persists the AI chat memory for a single user session. Rather than storing
every message individually in a separate table, the entity uses a hybrid memory model: a
\texttt{summary} text column holds a compressed rolling summary of older turns, updated by a
separate LLM call when the recent buffer overflows, and a \texttt{recentMessages} column stores
a JSON-serialised list of the most recent turns as plain text (compatible with PostgreSQL
\texttt{TEXT} and H2 for testing). The entity's primary key is a client-generated UUID string,
so no database sequence is required.

% TODO: Insert UML class diagram as Figure 3.1
% \begin{figure}[htbp]
%     \centering
%     \includegraphics[width=\textwidth]{images/fig_3_1_uml-class-diagram.png}
%     \caption{UML class diagram of the StudySpace domain model}
%     \label{fig:uml-class-diagram}
% \end{figure}

% ----------------------------------------------------------------
\section{Technology Analysis}
\label{sec:technologies}

This section justifies each major technology choice by comparing the available alternatives against
the requirements established in Chapter~\ref{ch:analysis}. The goal was to select a stack that is
well suited for the required features, has strong community support and ecosystem maturity, and is
practical to use within the scope of a bachelor thesis prototype.

% ----------------------------------------------------------------
\subsection{Frontend Framework}
\label{subsec:frontend-tech}

FR-01 through FR-04 require a responsive, component-rich user interface that manages complex
shared state across a course-material viewer, a chat panel, and an AI assistant panel rendered
simultaneously on the same screen.

\textbf{React~19}~\cite{React} is a declarative UI library maintained by Meta. Its component model
encourages decomposing the interface into small, reusable pieces, each responsible for a single
part of the view. The unidirectional data flow makes state changes predictable and easier to trace
during debugging. React's large ecosystem provides ready-made solutions for drag-and-drop panels,
rich-text editing, and real-time subscriptions, all of which are needed for StudySpace.

\textbf{Vue.js} is a progressive framework with a smaller bundle size and a gentler learning
curve. For StudySpace, however, the Vue ecosystem is shallower for complex, enterprise-grade
components, and the TypeScript integration, while improving, is less seamless than in React.

\textbf{Angular} is a full-featured framework backed by Google that includes dependency injection,
a router, and a form library out of the box. The overhead of Angular's concepts, decorators, and
module system exceeds what is justified for a single-team prototype of this scope.

React was selected because its component model aligns directly with the StudySpace screen layout
(course panel, chat sidebar, AI panel as discrete components), its TypeScript support is
first-class, and its ecosystem provides the resizable panel, rich-text input, and WebSocket hook
primitives needed by FR-02 and FR-03 without significant additional integration effort.

React is used through the \textbf{Next.js~15}~\cite{Nextjs} framework, which provides file-based
routing via the App Router, server-side rendering where applicable, and built-in image
optimisation. Compared to a plain Vite or Create React App setup, Next.js removes boilerplate
routing configuration and makes the project structure self-documenting through the directory
layout. Styling is handled by \textbf{Tailwind CSS}~\cite{TailwindCSS} combined with
\textbf{Shadcn/ui}~\cite{ShadcnUI} and \textbf{Radix UI}, which provide accessible, unstyled
component primitives that Tailwind classes are applied to. State management across the application
uses \textbf{Redux Toolkit}~\cite{ReduxToolkit}, which offers predictable reducers and typed
async thunks for all API interactions.

% ----------------------------------------------------------------
\subsection{Backend Framework}
\label{subsec:backend-tech}

FR-05 requires robust authentication and role-based authorisation. FR-02 and FR-04 involve
complex transactional operations over relational data. These requirements favour a statically
typed language with mature transaction management and security libraries.

\textbf{Java~21 with Spring Boot~3}~\cite{Java21,SpringBoot} provides a statically typed,
object-oriented foundation with a mature ecosystem. Spring Boot's convention-over-configuration
approach applies sensible defaults automatically, embeds a Tomcat web server, and offers
first-class integration with Spring Security, Spring Data JPA, Spring WebSocket, and Springdoc
OpenAPI --- all of which are essential for StudySpace. Virtual Threads, introduced in Java~21,
improve the throughput of concurrent request handling without the complexity of reactive
programming.

\textbf{Node.js with NestJS or Express} excels in I/O-bound tasks due to its event-driven
architecture. However, its single-threaded event loop can be a bottleneck for the multi-step
transactional operations required by the Content Extension System (FR-02), and the ORM ecosystem
does not reach the integration depth and stability of Spring Data JPA.

\textbf{Python with Django or FastAPI} is excellent for rapid prototyping and AI integration.
Python's dynamic typing, however, makes maintaining large-scale backend codebases more challenging
without strict tooling, and its Global Interpreter Lock can become a bottleneck for highly
concurrent CPU-bound tasks.

Java~21 with Spring Boot~3 was selected because the controller-service-repository layered model
maps directly onto the StudySpace domain structure, Spring Security provides JWT-based
authentication and method-level authorisation out of the box (directly addressing FR-05), and the
mature transaction management ensures referential integrity across the complex clone and proposal
approval operations required by FR-02.

% ----------------------------------------------------------------
\subsection{Database and Vector Store}
\label{subsec:database-tech}

The application's data model is inherently relational: users own workspaces, workspaces belong
to courses, clone records track inheritance, and proposals reference both student materials and
instructor courses. FR-02 requires ACID transaction guarantees across multi-entity writes.

\textbf{PostgreSQL~17}~\cite{PostgreSQL} is a mature, open-source relational database management
system (RDBMS) that follows the SQL standard and supports complex queries, transactions, and
foreign key constraints. It integrates seamlessly with the Java backend through the JDBC driver,
and higher-level access is provided by Spring Data JPA with Hibernate as the ORM provider.

\textbf{MongoDB} (NoSQL document store) and \textbf{Cassandra} (wide-column store) are optimised
for unstructured or semi-structured data and horizontal scalability. However, they lack strict
schema enforcement and standard join operations. Using a document store would require significant
denormalisation, duplicating data to avoid application-level joins, which introduces the risk of
inconsistency when a Course Material is approved and must be reflected in both the workspace and
the course tables atomically.

PostgreSQL was selected because the StudySpace data model is fundamentally relational, ACID
transactions are required for the clone and approval workflows (FR-02), and foreign-key
constraints enforce referential integrity between proposals and their source materials. The AI
assistant (FR-04) does not require a dedicated vector store at this prototype stage because the
document text is extracted at query time and injected directly into the prompt; full vector
embeddings are identified as a future extension in section~\ref{sec:future-work}.

% ----------------------------------------------------------------
\subsection{LLM API Integration}
\label{subsec:ai-tech}

FR-04 requires an LLM API capable of answering questions grounded in a provided document context
within a single prompt, without fine-tuning.

\textbf{Google Gemini API (gemini-2.5-flash)} provides a large context window, competitive
quality on open-domain question answering, and an official Java SDK that integrates directly with
the Spring Boot application. The free tier is sufficient for prototype-level usage.

\textbf{OpenAI GPT-4o} exposes a compatible REST interface and achieves comparable quality on
question answering tasks. However, the Java SDK is less mature than the official Google GenAI
SDK, and the cost per token is higher for the same context window size at the time of writing.

The Google Gemini API was selected because the official Google GenAI Java SDK~\cite{GoogleGenAI}
integrates directly with Spring Boot without additional HTTP client configuration, the
\texttt{gemini-2.5-flash} model's context window comfortably accommodates the 10\,000-character
document extract required by FR-04, and the free tier capacity is sufficient for the prototype.

% ----------------------------------------------------------------
\subsection{Authentication and Authorisation}
\label{subsec:auth-tech}

FR-05 requires stateless authentication suitable for a REST API consumed by a single-page
application, with role-based restrictions enforced on the server side.

\textbf{JSON Web Tokens (JWT) with Spring Security}~\cite{SpringSecurity} enables stateless
authentication: after login, the server issues a signed token that the client stores and sends
with every subsequent request. The server validates the token signature without consulting a
session store, which simplifies horizontal scaling.

\textbf{Server-side sessions} with Spring Session require a shared session store (such as Redis)
across instances and add network latency to every request. For a stateless REST API served to a
browser-based SPA, this overhead is unnecessary.

\textbf{OAuth~2.0 only} (without local credentials) would prevent users without a Google or
GitHub account from registering, which is inappropriate for a faculty application where students
may not use third-party providers for academic accounts.

JWT with Spring Security was selected because it achieves stateless session management (directly
satisfying NFR-02), requires no external session store, and Spring Security's filter chain
integrates the JWT validation and role check into a single, reusable \texttt{OncePerRequestFilter}
implementation. Google OAuth~2.0 login is supported as an optional alternative authentication
path alongside the primary local credential flow.

% ----------------------------------------------------------------
\subsection{Real-Time Communication}
\label{subsec:websocket-tech}

The study-session feature requires real-time event delivery to all connected participants without
polling, as required by the real-time study-session feature.


\textbf{WebSocket with STOMP over SockJS}~\cite{SockJS} layers a publish-subscribe model on top
of the raw WebSocket connection. Spring Boot registers a broker endpoint and routes messages to
per-session topics, allowing the server to push participant-list updates and activity events to
all subscribers without the clients polling.

\textbf{Server-Sent Events (SSE)} provide a unidirectional push channel from server to client.
They lack the bidirectional subscription model needed for clients to send activity events (such as
a hand-raise) back to the server over the same connection.

\textbf{Long Polling} achieves real-time delivery but requires each client to maintain a pending
request and re-establish it after every response, producing significant per-client overhead on
the server.

WebSocket with STOMP over SockJS was selected because the bidirectional subscription model covers
both server-to-client pushes (participant list updates, activity events) and client-to-server
messages (hand raises) over a single connection, and Spring Boot's built-in STOMP message broker
requires no additional infrastructure.

% ----------------------------------------------------------------
\subsection{File Storage}
\label{subsec:storage-tech}

FR-01 and FR-02 require file upload and retrieval for Course Materials and Private Workspace
materials. The storage backend must be swappable between a local development mode and a
production cloud target.

\textbf{Local file storage} (writing files to a directory under the working directory) is zero-cost
and requires no external service in development. It is unsuitable for production because files are
lost when the container is recreated and cannot be shared across multiple server instances.

\textbf{Google Cloud Storage (GCS)} is a managed object store with high durability, direct URL
serving, and a Java client library. It requires a GCP project and credentials, which adds setup
overhead during development.

A \texttt{FileStorageService} interface with two concrete implementations was designed:
\texttt{LocalFileStorageService} for development and \texttt{GcsFileStorageService} for production.
The active implementation is selected by a Spring profile, keeping the storage backend transparent
to the rest of the application and satisfying NFR-04. The interface defines three methods:
\texttt{store}, \texttt{copy}, and \texttt{delete}, which are the only storage operations needed
by FR-01 and FR-02.

% ----------------------------------------------------------------
\section{Database Schema}
\label{sec:database-schema}

The relational schema follows directly from the domain model described in
section~\ref{sec:domain-model}. Each entity maps to a dedicated table, and all foreign-key
relationships are enforced at the database level.

The primary tables are \texttt{users}, \texttt{courses}, \texttt{course\_sections},
\texttt{course\_materials}, \texttt{course\_enrollments}, \texttt{student\_workspaces},
\texttt{workspace\_spaces}, \texttt{workspace\_sections}, \texttt{workspace\_materials},
\texttt{contribution\_proposals}, and \texttt{conversations}. The
\texttt{course\_enrollments} table carries a unique constraint on
\texttt{(course\_id, student\_id)} to prevent duplicate enrollment rows. The
\texttt{workspace\_materials} table carries the two boolean flags \texttt{isReference}
and \texttt{isHidden} that implement the copy-on-write and soft-delete semantics described
in section~\ref{sec:domain-model}.

The AI conversation memory is stored in the \texttt{conversations} table as two \texttt{TEXT}
columns: \texttt{summary} for the compressed long-term memory and \texttt{recent\_messages} for
the JSON-serialised short-term buffer. This design avoids a separate message table and keeps the
memory retrieval to a single row lookup by the client-generated UUID primary key. The current
prototype does not use vector embeddings; document text is extracted at query time and injected
into the prompt. A dedicated embedding column using \texttt{pgvector} is identified as a future
extension that would replace on-the-fly extraction with indexed similarity search.

% TODO: Insert ERD as Figure 3.2
% \begin{figure}[htbp]
%     \centering
%     \includegraphics[width=\textwidth]{images/fig_3_2_er-diagram.png}
%     \caption{Entity-relationship diagram of the StudySpace database schema}
%     \label{fig:er-diagram}
% \end{figure}

% ----------------------------------------------------------------
\section{System Architecture (Final)}
\label{sec:system-architecture-final}

StudySpace follows a three-tier architecture. The \emph{presentation layer} is a Next.js and
React application running in the browser; it communicates with the backend via the REST API for
most operations and via a persistent WebSocket (STOMP/SockJS) connection during active study
sessions. The \emph{business logic layer} is a Spring Boot application that handles authentication,
data access, real-time session broadcasting, and LLM API calls. The \emph{data layer} is a
PostgreSQL database that stores all persistent data.

All three tiers are defined in \texttt{docker-compose.yml} and run in the same Docker network,
so they can reach each other by service name without exposing internal ports to the host.

The backend is organised into the following packages, each with a single responsibility and
dependencies flowing only from controller to service to repository:

\begin{table}[htbp]
    \centering
    \caption{Backend package structure}
    \label{tab:backend-packages}
    \renewcommand{\arraystretch}{1.4}
    \begin{tabularx}{\textwidth}{|p{3cm}|X|}
        \hline
        \textbf{Package} & \textbf{Responsibility} \\
        \hline
        \texttt{controller} & Defines REST endpoints, maps HTTP requests to service calls,
                              and validates incoming request bodies. \\
        \hline
        \texttt{service} & Contains all business logic. Services are called by controllers
                          and may call one or more repositories. Transactions are managed
                          at this layer. \\
        \hline
        \texttt{repository} & Spring Data JPA repositories, providing CRUD operations and
                              custom JPQL queries for complex data access. \\
        \hline
        \texttt{entity} & JPA entity classes mapped to database tables. \\
        \hline
        \texttt{dto} & Data Transfer Objects used as the API contract, decoupling the
                      API response shape from the internal entity structure. \\
        \hline
        \texttt{security} & JWT filter, \texttt{UserDetailsService} implementation, and
                           Spring Security configuration. \\
        \hline
        \texttt{config} & Application-level configuration beans: CORS policy, Swagger
                         setup, and WebSocket registration. \\
        \hline
    \end{tabularx}
\end{table}

% TODO: Insert final architecture diagram as Figure 3.3
% \begin{figure}[htbp]
%     \centering
%     \includegraphics[width=0.9\textwidth]{images/fig_3_3_system-architecture-final.png}
%     \caption{Final system architecture of StudySpace}
%     \label{fig:system-architecture-final}
% \end{figure}

% ----------------------------------------------------------------
\section{API Design}
\label{sec:api-design}

The backend exposes a RESTful API following standard HTTP conventions. Resources are identified
by URL paths, and operations are expressed through HTTP methods (\texttt{GET}, \texttt{POST},
\texttt{PUT}, \texttt{DELETE}). All protected endpoints require a valid JWT token in the
\texttt{Authorization: Bearer <token>} header.

Authentication follows a four-step flow. The client sends credentials to
\texttt{POST /api/auth/login}. The server validates the credentials and returns a signed JWT
token. The client stores the token in memory and includes it in the \texttt{Authorization}
header of all subsequent requests. The server's \texttt{JwtAuthenticationFilter} validates the
token on each request before passing control to the controller.

The full API is documented through the Swagger UI available at \texttt{/swagger-ui/index.html}
when the backend is running, auto-generated by Springdoc OpenAPI from controller annotations.

% ----------------------------------------------------------------
\section{UI and Wireframe Design}
\label{sec:ui-design}

The interface is role-aware: students and Instructors see different navigation options and action
buttons depending on their role. The application uses a sidebar navigation layout with the main
content area to the right.

\subsection{Navigation and Layout}
\label{subsec:navigation}

The sidebar items differ by role. Students see links to their enrolled courses, their Private
Workspaces, and their study sessions. Instructors additionally see course management, enrollment
management, and Merge Proposal review links. The role check is performed client-side from the
JWT payload, but all data-sensitive actions are re-validated server-side.

\subsection{Course and Workspace Views}
\label{subsec:course-workspace-views}

The course list page shows all courses the user is enrolled in or manages. Selecting a workspace
opens the workspace viewer, which renders the Course Material alongside a contextual chat panel.
Students can use the chat panel to discuss the course content and attach a Contextual Anchor to
a specific Course Material using the \texttt{@}-mention system.

% TODO: Insert workspace view wireframe as Figure 3.4
% \begin{figure}[htbp]
%     \centering
%     \includegraphics[width=0.9\textwidth]{images/fig_3_4_workspace-view.png}
%     \caption{Workspace viewer wireframe: Course Material list with contextual chat panel}
%     \label{fig:workspace-view}
% \end{figure}

The workspace viewer places the material list on the left and the contextual chat panel on the
right. This layout satisfies the FR-03 use case \texttt{Post contextual message} by keeping the
source material and the discussion visible simultaneously. The panel is resizable so students can
balance reading space against chat space according to their preference.

\subsection{Private Workspace and Merge Proposal Views}
\label{subsec:clone-workspace}

When a student clones a course, they receive a Private Workspace that they can edit freely. The
cloned workspace view mirrors the course structure and adds upload controls for new materials.
The Merge Proposal submission dialog allows the student to select materials and choose a target
section, optionally typing a proposed new section name if none of the existing sections is
appropriate.

% TODO: Insert merge proposal wireframe as Figure 3.5
% \begin{figure}[htbp]
%     \centering
%     \includegraphics[width=0.9\textwidth]{images/fig_3_5_merge-proposal.png}
%     \caption{Merge Proposal submission dialog wireframe}
%     \label{fig:merge-proposal}
% \end{figure}

The Merge Proposal dialog reflects the FR-02 use case \texttt{Submit Merge Proposal}: the student
selects one or more materials, picks a target section or proposes a new one, and submits.
The design keeps the form minimal to reduce friction and mirrors the Git-inspired workflow that
technically experienced students will recognise.

\subsection{AI Assistant Panel}
\label{subsec:ai-panel}

The AI assistant is accessible through a collapsible side panel within the workspace viewer. The
student types a question in natural language, and the system sends the question together with
the content of the currently open Course Material to the LLM API. The response is displayed in
a conversational format with a typewriter animation. The panel is designed not to obscure the
Course Material so the student can consult both simultaneously.

% TODO: Insert AI panel wireframe as Figure 3.6
% \begin{figure}[htbp]
%     \centering
%     \includegraphics[width=0.85\textwidth]{images/fig_3_6_ai-panel.png}
%     \caption{AI assistant panel wireframe: question input and response display}
%     \label{fig:ai-panel}
% \end{figure}

The collapsible design satisfies FR-04 by making the AI assistant available at any point during
material review without forcing the student to navigate to a separate page, preserving the reading
context while the answer is displayed.
